% ===========================================================================
% 07-Appendix.tex -- the apparatus chapter.
%
% main.tex issues \appendix before \include-ing this file, so the \chapter
% below numbers as A and its sections as A.1 ... A.6. Every \label here is
% referenced from the body; none may be renamed.
%
% \framedecl PRECEDES each \section deliberately, exactly as \beatsection
% orders them. Declared after the heading it would write the PREVIOUS
% section's frame into the mark that becomes \firstmark on this section's own
% opening page -- the bug documented in preamble.tex section 5.
%
% NO claim environment in this file, deliberately. Every claim block in the
% thesis sits in the Investigation chapter. An appendix records apparatus; it
% reaches no verdict the body has not.
% ===========================================================================

\chapter{Appendix}
\label{sec:appendix}

\noindent
This appendix holds the apparatus the body leans on but does not stop to
justify. Each section states its verdict in its heading. None is a new result;
several are defects that changed one.

% ===========================================================================
\framedecl{R}
\section{The bespoke estimators recover a planted answer, in both directions}
\label{sec:app-planted}

\begin{question}
Three measurements rest on estimators built for the purpose, and a reader has no
prior reason to trust them. The instrument is a world with a known answer
planted in it, run twice, with the effect and without. All three recover what
was planted, and nothing when nothing was.
\end{question}

\noindent
The weight falls on the null direction.\gloss{planted world}{A synthetic run
with the answer built in, no data, no network, one fixed seed.}

% ===========================================================================
\begin{table}[htbp]
\begin{threeparttable}
\caption[Planted-truth recovery]{\textbf{All three estimators recover planted
answers in both directions.} Every row is a synthetic planted-world check, and
none of these numbers comes from the evaluation run or is stored in a result
artifact. The $\Tcoef$ and
$\kappa$-structure rows reproduce from
\texttt{variance\_decomposition.py -{}-selftest}, the channel-gate rows from
\texttt{channel\_gate.py -{}-selftest}. The $\kappa$-structure check must
separate two worlds, within-line consistency exceeding across-line where every
cell line bends what it is given ($+1.17$), the excess vanishing where $\kappa$
is specific to each drug-and-line pairing with no line-level direction
($-0.00$). The lower row of each pair is the null world, the one that rules out
an instrument reporting structure wherever it looks.}
\label{tab:planted}
\small
\begin{tabularx}{\linewidth}{@{}>{\raggedright\arraybackslash}X r@{}}
\toprule
\thead{Planted truth} & \thead{Recovered} \\
\midrule
\multicolumn{2}{@{}l}{\itshape transfer coefficient $\Tcoef$} \\
\quad $\Tcoef = 1.00$ (fully shared) & $0.999$ \\
\quad $\Tcoef = 0.70$ & $0.703$ \\
\quad $\Tcoef = 0.40$ & $0.404$ \\
\addlinespace
\multicolumn{2}{@{}l}{\itshape channel gate, over 24 worlds} \\
\quad a channel that genuinely predicts & $+0.397$ (sd $0.029$) \\
\quad a channel carrying nothing & $-0.016$ (sd $0.091$) \\
\addlinespace
\multicolumn{2}{@{}l}{\itshape $\kappa$-structure} \\
\quad a consistent per-line modulation & $+1.173$ \\
\quad an idiosyncratic interaction & $-0.002$ \\
\bottomrule
\end{tabularx}
\end{threeparttable}
\end{table}

\paragraph{What the transfer-coefficient rows do not show} They are a data-free
unit test whose planted world has cell line as its only unshared axis, so
recovery shows the estimator inverting its own generative model and no more. The
canonical run, computed with \texttt{same\_dose\_only=false},
\texttt{loo\_generic=false} and \texttt{repro\_thr=0.2}, lets cell line, dose,
well and estimator noise all enter one number. On synthetic data calibrated to
its gene dimension, cell counts and empirical noise, the estimator returns
$1.001$, $0.747$ and $0.494$ for planted $\Tcoef$ of $1.00$, $0.75$ and $0.50$.

\begin{scopelimit}[carried into \cref{sec:res-transfer}]
Neither the planted rows nor the calibrated repeat licenses reading
$1 - \Tcoef$ as a drug $\times$ cell-line interaction share. The quantity
measures the fraction not shared across the compared contexts, and in the
canonical run those contexts differ in more than cell line.
\end{scopelimit}

\paragraph{Two of these checks changed the instrument} The count-matched null
originally drew random partners from a pool \emph{excluding} the channel's own
selections, which depletes the null of the partners the channel chose and
anti-correlates the two arms. And the useless-channel gap has a standard
deviation of roughly $0.09$ on one small world, so the original validation,
asserting that a single draw fell within $0.10$, tested luck; it now averages 24
worlds. That same variance is why the channel results are gated on a clustered
interval and not on a point estimate.

\paragraph{A third defect was found after the canonical run} The \emph{retrieval}
pool, whose partners' mean residual forms the prediction, was every retained
condition in the cell line except the target's own drug, so a condition held out
of training could enter the prediction for a held-out drug. That is not what
\cref{sec:res-scope} says the gate was intended to do. The code now restricts
the pool to the split's training conditions, logs its size, and never falls back
silently.

The training subset is
\needsaudit{2{,}523} of a retained total logged as
\needsaudit[the split tally sums differently]{4{,}131} conditions, so the old
pool was around $64\%$ larger than it was entitled to be. What makes it a defect
is that a held-out condition could enter at all, which holds at any tally. Every
channel-gate figure in the main text is from the corrected
re-run\corrmark{The channel gate's retrieval pool, repaired and re-run after the
canonical run. \cref{sec:res-scope}}; the planted worlds are unaffected. The
repair moved the channels in no common direction. Target fell, mechanism rose,
chemistry roughly halved, so the pre-fix figures are not quoted alongside them.

% ===========================================================================
% PLACEMENT ONLY. The global \section reservation is 7\baselineskip, sized in
% preamble.tex section 4 for a ONE-LINE heading plus the opener's three-line
% minimum. This title sets on two lines, so 7 lines cleared the test and then
% left the heading, the beat strip and exactly two lines of the lede at the foot
% of the page, with the rest of the lede and the whole VERDICT banner overleaf.
% 15 lines is the two-line heading plus the entire opener, which is how every
% other section in this appendix sets.
\needspace{15\baselineskip}%
\framedecl{E}
\section{The absent-gene convention does not change the answer}
\label{sec:app-convention}

\begin{question}
A predicted cell sentence lists only the genes it expresses, so some panel genes
reach the scorer with no rank and the scorer invents one. Does the invented rank
decide the result? Two conventions were carried side by side. The choice does
not matter, and it never reaches the metric the canonical results rest on.
\end{question}

\noindent
\texttt{worst} assigns absent genes the last rank $G$; \texttt{francesca}
assigns them a fixed middle rank $G/2$. These are the extremes of the plausible
range, so a convention that mattered anywhere would show up between them.

It does not. Spike-in $\DEdr$ reads $0.952$ under the first and $0.954$ under
the second, and no baseline or positive control scored under both moves by more
than $0.01$. That comparison was a one-off diagnostic, never written to an
artifact, so it is a check that was run and not a reproducible result.

\paragraph{How far the convention reaches} It is an input to the rank-based
scores $\DEdr$, $\paneltau$ and \texttt{nir\_rank}, which fills absent genes at
the last rank like the rest.\seealso{\cref{sec:methods-metrics} for why
\texttt{nir\_rank} never entered the graded family} It does not reach the $\NIR$
figures the canonical results rest on, scored by instruments that never receive
the fill. All body figures use \texttt{worst}.

% ===========================================================================
\framedecl{E}
\section{The swap-distance sweep bounds a large effect, and only in one column}
\label{sec:app-sweep}

\begin{question}
If the scramble control swaps in a drug that behaves like the real one, the null
it produces says nothing. The instrument is an earlier, continuous version of
that control, measuring the scramble gap against swap distance. It bounds a
large effect; it does not establish a null, and only one of its two columns
carries that bound.
\end{question}

\noindent
None of these numbers is canonical.\gloss{swap distance}{How far the swapped-in
drug's response lies from that of the drug it replaced.} This is a
separate, earlier run. Its columns are the single-cell and
optimal-transport checkpoints, not the residual arm; it was scored on the
unseen-drug tier of a different evaluation set; and its swap partners were not
restricted to plate, whereas the canonical run draws within plate. Both
artifacts, \texttt{scramble\_sweep\_single.json} and
\texttt{scramble\_sweep\_T2.json}, record \texttt{max\_new\_tokens=1200},
against the $1400$-token budget of the later optimal-transport result.

% ===========================================================================
\begin{table}[htbp]
\begin{threeparttable}
\caption[Swap-distance sweep]{\textbf{No bin in either model has an interval
excluding zero in this historical $1200$-token run, and the trend correlations
are $\approx 0$.} 200 (drug, swap) pairs over 20 cell lines, $n=40$ per bin. The
single-cell column bounds a large distance effect here. The optimal-transport
column does not carry into the canonical result, where the same checkpoint
clears zero at $1400$ tokens under the within-plate stratified test
(\cref{sec:res-encoding}).}
\label{tab:sweep}
\small
\begin{tabularx}{\linewidth}{@{}%
  >{\raggedright\arraybackslash}X r r@{}}
\toprule
\thead{Bin (swap distance)} & {\thead{single-cell model}} &
  {\thead{optimal-transport model}} \\
\midrule
0, nearest (2.75--4.03) & $-0.046$ & $-0.001$ \\
1 (4.03--4.60) & $+0.060$ & $+0.031$ \\
2 (4.60--4.93) & $-0.006$ & $+0.017$ \\
3 (4.93--5.32) & $-0.004$ & $+0.043$ \\
4, farthest (5.32--6.25) & $-0.019$ & $-0.005$ \\
\addlinespace
trend, corr(distance, $\gap$) & $+0.051$ & $+0.012$ \\
\bottomrule
\end{tabularx}
\end{threeparttable}
\end{table}

% ===========================================================================
\begin{figure}[htbp]
\centering
\begin{tikzpicture}[x=21mm, y=1mm, font=\small]
  % The DOT carries the value, always. Four of the five gaps sit within a few
  % thousandths of zero, so their labels are lifted to one clear row and tied
  % back with a leader. Nothing is repositioned that could be misread as a
  % different number. There are no numeric axis ticks; each strip carries the
  % one reference line its own quantity has, zero for the gap above and chance
  % for the $\NIR$ below, and the lower strip is a DIFFERENT QUANTITY, which is
  % what the divider between the two strips is there to say.
  \def\ur{180}                     % upper strip, mm per unit of the gap
  \def\lz{-52}\def\ls{63}          % lower strip, origin and mm per unit
  \def\lchance{\lz+(0.5-0.326)*\ls}% -41.04mm: $\NIR = 0.50$ on the lower strip
  \newcommand{\liftedgap}[2]{%
    \fill[ink] (#1,{#2*\ur}) circle (1.5pt);
    \draw[rulegrey, line width=0.3pt] (#1,{#2*\ur}) -- (#1+0.08,-14);
    \node[anchor=west, font=\scriptsize, text=ink] at (#1+0.10,-14) {$#2$};}
  % The value label carries a white backdrop because the chance rule runs
  % under this strip: the nearest value, 0.497, sits ON that rule, and a dotted
  % line through its digits is the one way this label could be misread.
  \newcommand{\nirdot}[2]{%
    \fill[quietink] (#1,{\lz+(#2-0.326)*\ls}) circle (1.3pt);
    \node[anchor=west, font=\scriptsize, text=quietink, fill=white, inner sep=1pt]
      at (#1+0.06,{\lz+(#2-0.326)*\ls}) {$#2$};}

  % === upper strip, the scramble gap in the single-cell column ==============
  \node[anchor=west, font=\scriptsize\bfseries, text=accent]
    at (-0.42,29) {$\gap$, single-cell column};
  \draw[ink, line width=0.5pt, densely dotted] (-0.42,0) -- (4.40,0);
  \node[anchor=west, font=\scriptsize\itshape, text=ink] at (4.42,0) {zero};

  % the one bin for which an interval is quoted, drawn as an interval
  \draw[accent!55, line width=1.6pt] (1,{-0.002*\ur}) -- (1,{0.124*\ur});
  \fill[accent] (1,{0.060*\ur}) circle (1.6pt);
  \node[anchor=west, font=\scriptsize, text=accent]
    at (1.08,{0.060*\ur}) {$+0.060$};

  \liftedgap{0}{-0.046}
  \liftedgap{2}{-0.006}
  \liftedgap{3}{-0.004}
  \liftedgap{4}{-0.019}

  % === the divider IS the argument, two quantities, two scales, one x =======
  \draw[rulegrey, line width=0.4pt] (-0.42,-20) -- (4.40,-20);

  % === lower strip, the model's own score, which drifts =====================
  \node[anchor=north west, font=\scriptsize\bfseries, text=quietink]
    at (-0.42,-23) {the model's own $\NIR$ (expression frame) in the same bins,
      on its own scale};
  % The chance line at 0.50 is drawn on EVERY $\NIR$ axis, always
  % (figs/style_v2.py): omitting it, or implying it through the extent of the
  % marks, lets a reader infer a scale that does not exist. Three of these five
  % values sit below it and two above, and nothing else in the drawing says so.
  % Same dotted ink as the zero rule above, because the two are the same kind of
  % thing -- each strip's own single reference -- and drawn BEFORE the marks so
  % the data reads on top of it.
  \draw[ink, line width=0.5pt, densely dotted]
    (-0.42,{\lchance}) -- (4.40,{\lchance});
  \node[anchor=west, font=\scriptsize\itshape, text=ink]
    at (4.42,{\lchance}) {chance};
  \nirdot{0}{0.497}
  \nirdot{1}{0.612}
  \nirdot{2}{0.523}
  \nirdot{3}{0.401}
  \nirdot{4}{0.326}

  % === the one axis the two strips share ====================================
  \foreach \i/\lab in {0/nearest, 1/{}, 2/{}, 3/{}, 4/farthest}{
    \node[anchor=north, font=\scriptsize, text=quietink, align=center]
      at (\i,-58) {\i\\\lab};}
\end{tikzpicture}
\caption[Swap distance and the scramble gap]{\textbf{The one borderline value
sits in an inner bin, and the drift that could have faked a trend is in a
quantity the gap cancels.} Upper strip, the scramble gap in the single-cell
column across the five swap-distance bins of \cref{tab:sweep}; the bar marks the
only bin for which an interval is quoted, \ci{+0.060}{-0.002}{+0.124}. A genuine
dissimilarity effect would be monotone and peak at the right-hand edge, and the
one value approaching zero from above is the second bin of five. Lower strip,
the model's own $\NIR$ in those bins ($0.497$, $0.612$, $0.523$, $0.401$,
$0.326$), falling overall but not monotonically, because drugs far from
everything populate the far bins. The strips are different quantities on
different scales; $\gap$ uses the same drug on both arms, so that drift cancels.
Neither strip is canonical (\texttt{max\_new\_tokens=1200}, $n=40$ per bin).}
\label{fig:app-sweep-bins}
\end{figure}

\paragraph{What the sweep can and cannot support} $n=40$
per bin gives half-widths of $0.06$--$0.10$ in the single-cell column and
$0.03$--$0.06$ in the optimal-transport column, so the sweep excludes a large
effect without establishing a
null.\seealso{\cref{sec:methods-controls} for the stratified scramble
that replaced this design} The later budget-matched result withdraws the
optimal-transport column as evidence of a null, that verdict being
generation-budget dependent and non-null at $1400$ tokens
(\cref{sec:res-encoding}).\corrmark{The optimal-transport column of this sweep,
withdrawn as evidence of a null. \cref{sec:res-encoding}}

\begin{scopelimit}[carried into \cref{sec:res-blind}]
Only the single-cell column remains a bound from this earlier design, and it is
a bound and not a null. At this sample size it excludes a large dissimilarity
effect and leaves a small one open. The optimal-transport column has been
withdrawn as evidence of a null, so nothing in this section rests on it.
\end{scopelimit}

% ===========================================================================
\framedecl{-}
\section{Undertraining is excluded as a trivial explanation, not as a general one}
\label{sec:app-training}

\begin{question}
A negative result invites the objection that the model was simply not trained
enough, and that objection deserves an answer from the logs. Two instruments
give one, what the loss was doing while the learning rate was still large, and
what a much longer schedule does to held-out cross-entropy. Under this schedule,
longer optimisation does not rescue it.
\end{question}

% ===========================================================================
% MOVED HERE from the end of A.4 -- PLACEMENT ONLY. Not one character of the
% figure, its caption or its numbers was changed; only the point of declaration.
%
% It is declared AHEAD of tab:trainconfig deliberately. Both are section-opening
% material, but the figure has a \cref in this section and the table does not:
% tab:trainconfig is a lookup cited only from chapters 4 and 5, so it is free to
% take the later float slot while the figure takes the one nearer its reference.
%
% MEASURED, in this order:
%   declared at the end of A.4 (as it was): float lands p.118, two pages past
%     its \cref on p.116, and the room it leaves under itself on p.118 fits the
%     two \Large lines of the A.5 heading and nothing else -- the heading is
%     stranded at the foot of the page with zero lines below it.
%   declared next to its \cref but still after the table: float still lands
%     p.118. The table, declared earlier, holds the only top slot on p.117.
%   \needspace{10\baselineskip} on the A.5 heading: BYTE-IDENTICAL PDF. It
%     cannot work from the heading's side -- \needspace tests
%     \pagegoal-\pagetotal in the main vertical list, before the deferred float
%     has been contributed to the page, so it reads a full page of room and
%     passes. The float is the cause, so the float is what moves.
%   declared here: float lands p.117, one page from its \cref, and the A.5
%     heading is no longer orphaned. Cost: one page, and p.117 sets the figure
%     with two lines of text under it.
\begin{figure}[htbp]
\begin{fullwidth}
\centering
% drawn at 174mm by thesis_v2/figs/training.py; no width= key, so it is placed
% at scale 1.0 and its type prints at the point sizes it was set in.
\includegraphics{fig-training}
\caption[The short-run arms flatten; a longer schedule
overfits]{\textbf{The canonical short-run arms flatten, while a
separate longer residual schedule overfits on held-out cross-entropy.}
(a) De-averaged training loss against fraction of the nominal one-epoch horizon,
three target constructions, each on \emph{its own} $y$-axis: absolute losses are
not comparable between arms because their target strings differ, so no shared
scale is offered. The residual and optimal-transport arms complete their epoch;
consensus was scheduled for one but terminated by the wall-clock at $61.5\%$,
where its line stops. Three superseded residual builds are omitted, as is the
single-cell arm, whose training log was not retained. Every arm runs a cosine
schedule annealing to $\approx 0$; the shaded band marks $50$--$70\%$ of the
horizon, where the residual arm's learning rate is still $22$--$52\%$ of its
post-warmup peak of $10^{-5}$ and its de-averaged loss travels only from $1.859$
to $1.854$. (b) Epoch-mean losses for the ten-epoch residual schedule (nominal
horizon $92{,}318$; $92{,}310$ optimizer updates). The ringed point is that
run's best validation loss, at epoch 1. The dashed line is the canonical
one-epoch run's validation loss on the same $3{,}600$-example file, $1.8986$,
which the ten-epoch run does not reach at any epoch. Training loss is above
validation at epoch 1 and below it at epoch 2, so the crossing falls inside the
shaded epoch; where inside it is not logged, and no point is drawn for it. Two
things this figure does not establish: that all possible learning-rate schedules
are ruled out, and that $\NIR$ worsens with additional epochs --- no long-run
checkpoint was scored with $\NIR$.}
\label{fig:training}
\end{fullwidth}
\end{figure}

% ===========================================================================
\begin{table}[htbp]
\begin{threeparttable}
\caption[Training configuration]{\textbf{One base checkpoint, one
hyperparameter set, five arms differing only in target.} Learning rate
$10^{-5}$, batch size 1, gradient accumulation 16, maximum length 8192, bf16,
weight decay 0.01, warmup ratio 0.03 throughout. The consensus arm stopped early
on a flat loss, so it is the $\varepsilon$-ladder's weakest-evidenced endpoint,
and at the residual arm's generation budget against a neutral partner the
optimal-transport endpoint is not null either (\cref{sec:res-encoding}); outside
\cref{tab:epsilon} the ladder's negative rests on the single-cell endpoint
alone. These five arms (\texttt{single\_cell}, \texttt{consensus},
\texttt{ot\_T2}, \texttt{residual}, \texttt{residual\_holdout}) are also the
complete family over which the identity test of \cref{sec:res-mechanism} takes
its multiplicity correction, five arms by up to seven depths and $26$ measured
tests, the family behind the $q$-values quoted there. One qualification attaches
to the residual row. Its step count is the rebuilt target build's, on which the
canonical behavioural results are scored, while the identity test of that arm
ran on the earlier target build.}
\label{tab:trainconfig}
\footnotesize
\begin{tabularx}{\linewidth}{@{}%
  >{\raggedright\arraybackslash}X
  >{\raggedright\arraybackslash}X
  l
  >{\raggedright\arraybackslash}X@{}}
\toprule
\thead{Arm} & \thead{Target} & \thead{Training} & \thead{Note} \\
\midrule
single cell & the observed treated cell & 1 epoch & \\
consensus & group pseudobulk & \textbf{$\approx 60\%$ of an epoch}
  & tier-1 $\NIR$ unscoreable \\
optimal transport (T2) & barycentric, $\varepsilon{=}0.5$
  & 1 epoch, 23{,}842 steps & \\
residual & signed residual & 1 epoch, 9{,}231 steps
  & rebuilt build's step count \\
residual (holdout) & as above, 738 withheld & 1 epoch & \\
\bottomrule
\end{tabularx}
\end{threeparttable}
\end{table}

\begin{scopelimit}[carried into \cref{sec:res-mechanism}]
The residual arm's two measurements are not on the same weights. The step count
recorded above is the rebuilt target build's, which is the build the canonical
behavioural results are scored on, whereas the identity test of that arm ran on
the earlier target build. A behavioural result for that arm and an activation
probe of it therefore come from two different builds, and this qualification
travels with any statement that pairs them.
\end{scopelimit}

\paragraph{How panel (a) must be read} Two properties of the one-epoch logs behind
\cref{fig:training} decide it. The trainer records a
\emph{cumulative} epoch mean, not an instantaneous loss, and that flattens by
arithmetic whatever the model is doing, so every series is de-averaged to its
per-window value before plotting. And the cosine schedule anneals the learning
rate to $\approx 0$ by the end of the epoch, $2.96 \times 10^{-10}$ at the last
logged step against a post-warmup peak of $10^{-5}$, so the final third of each
curve is flat partly because the step size has gone to zero, and is not
independent evidence of anything.

What is load-bearing is the band between $50\%$ and $70\%$ of the epoch, where
the learning rate is still $22$--$52\%$ of peak, enough to move the weights, and
the residual arm's loss has already flattened, travelling only from $1.859$ to
$1.854$. Its final held-out loss, on a matched validation file of $3{,}600$
examples, is $1.8986$ against a training loss of $1.8854$.

\paragraph{Panel (b) tests ordinary longer optimisation directly} Both residual
logs agree on base model and tokenizer,
training and validation files and counts, seed, effective batch, and peak logged
learning rate; weight decay and AdamW betas are not logged and are not claimed
as verified matches (\texttt{training\_adequacy\_}\allowbreak\texttt{long.json}). The ten-epoch run
has a nominal horizon of $92{,}318$ but executes $92{,}310$ optimizer updates,
accumulation resetting at every epoch boundary so that each epoch contributes
$\lfloor147{,}710/16\rfloor=9{,}231$. It is a separate longer schedule, not a
continuation; changing \texttt{num\_epochs} stretches the cosine schedule across
all ten epochs.

The result runs the wrong way for the objection. Validation loss is best at
epoch 1 ($1.9199$), worsens to $2.0559$ and $2.1708$ at epochs 2 and 3, and ends
at $2.2568$, while training loss falls from $1.9456$ to $1.4485$. The canonical
one-epoch run remains better on the same validation file ($1.8986$).

\begin{scopelimit}[carried into \cref{sec:limitations}]
That result does not show that every possible schedule would fail, and no
long-run checkpoint was scored with $\NIR$, so it says nothing about whether
$\NIR$ itself worsens with additional epochs. Together with the one-epoch
traces, it excludes undertraining as a \emph{trivial} explanation, not as a
general one.
\end{scopelimit}

% ===========================================================================
% PLACEMENT ONLY -- MOVED, not edited. This float used to be declared here, at
% the very end of A.4, some 40 lines after the \cref that introduces it. Two
% measured consequences, both gone now that it is declared next to its first
% reference (see the block above \paragraph{Panel (b)}):
%   1. It could not be placed until the builder was already on p.118, two pages
%      past the \cref on p.116.
%   2. Placed at the top of p.118 it left just enough room under itself for the
%      two \Large lines of the A.5 heading and nothing else, stranding that
%      heading at the foot of the page. \needspace cannot repair this from the
%      heading's side: it tests \pagegoal-\pagetotal in the main vertical list,
%      before the deferred float has been contributed to the page, so it reads a
%      full page of room and passes. VERIFIED: \needspace{10\baselineskip} on
%      the A.5 heading produced a byte-identical PDF. The float is the cause and
%      the float is what has to move.
% Neither the figure, its caption, nor one number in either was touched.

% ===========================================================================
\framedecl{-}
\section{Three identity-test designs failed, and each failure is a constraint on the fourth}
\label{sec:app-identity}

\begin{question}
The substitution test the body reports is the fourth design, and the three
before it each produced a result this thesis once quoted as a finding. The
instrument here is the algebra of each failure, what the control arm was
actually doing.
\end{question}

\noindent
All three failures below are failures of the control arm, not of the readout.
\Cref{fig:slab} draws the geometry each of them misses, and its three insets are
the three designs in the order they are argued here.

% ===========================================================================
% slab.tex -- "represented but not consulted", drawn.
% \input by Sections/04b-representation.tex inside \beatsection sec:res-used.
% Compiles against thesis_v2/preamble.tex and nothing else: no \includegraphics,
% no external data, no TikZ library beyond the ones preamble.tex already loads
% (this file uses arrows.meta, calc and decorations.pathreplacing).
%
% EVERY NUMBER DRAWN, AND WHERE IT COMES FROM
% -------------------------------------------
% Source: RESULTS_cluster/workspace_probe.json, key layers.<L>.variance_share,
% over the seven layers the probe ran (2, 4, 6, 8, 9, 12, 16). Config there:
% n_drugs 12, n_per_drug 40, n_dims 10, chance 0.08333. Sibling ledger of the
% exact values drawn: thesis_v2/figs/slab.json.
%
%   quantity            per-layer values (2,4,6,8,9,12,16)              mean
%   context        .701714 .738805 .756182 .749826 .741689 .838166 .838560  0.76642
%   cell_line      .579634 .641342 .665019 .667029 .654599 .567865 .492515  0.60971
%   drug_pure      .035747 .030703 .030347 .029881 .030620 .029807 .033109  0.03146
%
% The two WEDGE AREAS are the only scaled marks in the top panel:
%   context wedge   0.766420 of the disc -> 275.9113 deg
%   V_pure wedge    0.031459 of the disc ->  11.3252 deg
%   unshaded rest   0.202121 of the disc ->  72.7635 deg   (1 - the other two;
%                   the two subspaces are orthogonal, so the shares add and the
%                   remainder is what is left. It carries NO printed number.)
% Sum check: 275.9113 + 11.3252 + 72.7635 = 360.0000.
%
% AGREEMENT WITH THE TEXT (checked, not assumed):
%   04b-representation.tex:341  "stays flat at ~0.03 against ~0.6 for cell line"
%                               -> drawn 0.031 and quoted 0.610.               OK
%   04b-representation.tex:350  "cell line measures 15-22x the purified slab"
%                               -> per-layer ratios 14.9 16.2 19.0 20.9 21.4
%                                  21.9 22.3, i.e. 15-22x to two figures.      OK
%   04b-representation.tex:376  "flat at ~0.03, roughly 5% of the cell line's
%                               ~0.6"                                          OK
%   07-Appendix.tex:446-448     "roughly 99.5% of the control's energy ...
%                               cell-line directions at ~0.6 variance share"   OK
%   07-Appendix.tex:456-457     control delivers 5.3x the perturbation         OK
%   07-Appendix.tex:466         normaliser inflation "+2.75 nats"              OK
%   04b-representation.tex:260  "twelve centred class means span at most eleven
%                               directions"; top TEN kept (pure_drug_dims 10)  OK
% The context subspace's own share, 0.766, appears in NO sentence of the thesis;
% it is a new reading of the same JSON field, and it is labelled as the context
% subspace so it cannot be mistaken for the cell-line 0.61 the text quotes.
%
% INSET 1's ANGLE is a drawn quantity, not a decoration: a direction drawn
% isotropically in 2048 dimensions puts 10/2048 = 0.49% of its energy in a
% ten-dimensional slab, which is the appendix's "roughly 99.5%". An arrow at
% 86 deg to the slab axis has cos^2(86) = 0.0049 of its energy along that axis,
% so the arrow IS the number. Nothing else in the insets is to scale.
%
% LAYOUT NOTES, all forced by a build:
%   * fullwidth = textwidth 142mm + marginparwidth 28mm + marginparsep 4mm
%     = 174mm. Content is kept inside x in [0.00, 17.28]; a text-width block
%     starting at 12.00 must not exceed 5.20cm, because at 5.30cm the node's
%     inner sep pushed the picture 0.75pt over the measure and TeX said so.
%     (The "171mm" in preamble.tex's comment at \newenvironment{fullwidth} is
%     stale arithmetic; the geometry keys give 174mm.)
%   * the slab wedge is 11.3 deg wide and at radius 3cm is 5.9mm across at the
%     rim. Three class means and two arrows do not fit in it at true scale, so
%     the slab's interior is drawn ONCE, magnified, in the dashed panel, with
%     the two leader lines that make it a magnifier. The magnification is the
%     whole reason the size asymmetry can stay true: nothing was enlarged to
%     make a label fit.
%   * insets 2 and 3 repeat the leader-line device into their own small strips,
%     inset 1 does not. That is meaningful and not an inconsistency: failure
%     one happens OUTSIDE the slab and is legible at disc scale, failures two
%     and three happen inside it.
%   * the arm names sit UNDER mu_B and mu_C, not along the arrows. Along the
%     arrows they collided with each other in the wedge between them, because
%     the two displacements share an anchor; and mu_A had to be pushed left of
%     its dot, because directly above it is where Pi(mu_B - mu_A) begins.
% ===========================================================================
\begin{figure}[htbp]
\begin{fullwidth}
\centering
% --- the three wedge angles, computed above, written once --------------------
\newcommand{\slabHalf}{5.6626}     % half of 11.3252 deg  = 0.031459 of the disc
\newcommand{\slabCtxA}{5.6626}     % context wedge starts where the slab ends
\newcommand{\slabCtxB}{281.5739}   % 5.6626 + 275.9113
% --- one disc, drawn at whatever radius is asked for -------------------------
\newcommand{\slabdisc}[1]{%
  \filldraw[fill=panelbg, draw=rulegrey, line width=0.35pt]
    (0,0) -- (\slabCtxA:#1)
    arc[start angle=\slabCtxA, end angle=\slabCtxB, radius=#1] -- cycle;
  \filldraw[fill=accent, draw=accent, line width=0.3pt]
    (0,0) -- (-\slabHalf:#1)
    arc[start angle=-\slabHalf, end angle=\slabHalf, radius=#1] -- cycle;
  \draw[rulegrey, line width=0.5pt] (0,0) circle (#1);}

\begin{tikzpicture}[x=1cm, y=1cm, font=\footnotesize,
                    every node/.style={inner sep=1.6pt}]
  \tikzset{
    lab/.style     ={text=ink, font=\footnotesize},
    quiet/.style   ={text=quietink, font=\scriptsize},
    ttl/.style     ={text=accent, font=\scriptsize\bfseries\sffamily},
    lead/.style    ={draw=rulegrey, line width=0.3pt},
    swaparm/.style ={draw=accent, line width=1.0pt,
                     -{Latex[length=3.4pt,width=2.8pt]}},
    ctrlarm/.style ={draw=ink, line width=1.0pt,
                     -{Latex[length=3.4pt,width=2.8pt]}},
    dot/.style     ={circle, fill=ink, inner sep=0pt, minimum size=3.0pt},
    adot/.style    ={circle, fill=accent, inner sep=0pt, minimum size=3.4pt}}

% ###########################################################################
% # TOP BAND -- the stream, and the slab magnified
% ###########################################################################

  % ---------------- the disc ------------------------------------------------
  \begin{scope}[shift={(3.55,8.50)}]
    \slabdisc{3.0}

    \node[ttl, anchor=south] at (0,3.16)
      {THE 2048-DIMENSIONAL RESIDUAL STREAM};

    % the context subspace, labelled inside its own wedge
    \node[lab, align=center, text width=3.2cm] at (150:1.18)
      {\textbf{context subspace}\\ cell line, plate, dose};

    % the unshaded remainder carries no number
    \node[quiet, align=center, text width=2.1cm] at (318:1.95)
      {the remaining\\ directions};

    % the slab itself, too thin to hold a label
    \node[text=accent, font=\footnotesize, anchor=west]
      at (0:3.14) {$V_{\mathrm{pure}}$};

    % leaders to the magnified panel. Panel corners in OUTER coordinates are
    % (8.15,11.00) and (8.15,5.65); this scope's origin sits at (3.55,8.50).
    \draw[lead] (\slabHalf:3.0)  -- (4.60,2.50);
    \draw[lead] (-\slabHalf:3.0) -- (4.60,-2.85);
  \end{scope}

  % ---------------- the slab, magnified -------------------------------------
  \draw[rulegrey, line width=0.35pt, densely dashed, rounded corners=2pt]
    (8.15,5.65) rectangle (14.10,11.00);

  \node[ttl, anchor=north west] at (8.35,10.86) {THE SLAB, MAGNIFIED};
  \node[quiet, anchor=north west, align=left, text width=5.5cm] at (8.35,10.48)
    {at most ten of the $2048$ directions; $0.031$ of the stream's variance.
     The geometry below is schematic.};

  \coordinate (muA) at (9.40,7.90);
  \coordinate (muB) at ($(muA)+(30:2.05)$);
  \coordinate (muC) at ($(muA)+(-16:2.05)$);

  % equal norm, drawn as the one arc both heads land on
  \draw[rulegrey, line width=0.35pt, densely dashed]
    ($(muA)+(-16:2.05)$) arc[start angle=-16, end angle=30, radius=2.05];
  \node[quiet, anchor=west] at ($(muA)+(7:2.17)$) {identical norm};

  \draw[swaparm] (muA) -- (muB);
  \draw[ctrlarm] (muA) -- (muC);
  \node[lab, rotate=30,  anchor=south, inner sep=3pt]
    at ($(muA)!0.50!(muB)$) {$\Pi(\mu_B-\mu_A)$};
  \node[lab, rotate=-16, anchor=north, inner sep=3pt]
    at ($(muA)!0.50!(muC)$) {$\Pi(\mu_C-\mu_A)$};

  \node[adot] at (muA) {};
  \node[dot]  at (muB) {};
  \node[dot]  at (muC) {};
  \node[lab, anchor=east]       at ($(muA)+(-0.16,-0.06)$) {$\mu_A$};
  \node[lab, anchor=west]       at ($(muB)+(0.10,0.06)$)   {$\mu_B$};
  \node[lab, anchor=west]       at ($(muC)+(0.10,0.06)$)   {$\mu_C$};
  \node[quiet, anchor=west]     at ($(muB)+(0.10,-0.20)$)  {\textsc{swap} arm};
  \node[quiet, anchor=west]     at ($(muC)+(0.10,-0.18)$)  {\textsc{permuted} arm};

  \node[quiet, anchor=north west, align=left, text width=5.5cm] at (8.35,6.95)
    {Equal length is what makes \textsc{permuted} a matched control: same
     anchor, same construction, same norm, wrong drug.};

  % ---------------- the measured shares -------------------------------------
  \node[ttl, anchor=north west] at (14.30,10.86) {VARIANCE SHARE};
  \node[quiet, anchor=north west] at (14.30,10.50) {seven layers, mean};
  \draw[rulegrey, line width=0.4pt] (14.30,10.16) -- (17.25,10.16);
  \node[anchor=north west, text=ink, font=\scriptsize] at (14.30,10.06)
    {\begin{tabular}{@{}l@{\hspace{20pt}}r@{}}
       context & $0.766$\\[1.5pt]
       \textcolor{quietink}{cell line} & \textcolor{quietink}{$0.610$}\\[1.5pt]
       $V_{\mathrm{pure}}$ & \textcolor{accent}{$0.031$}\\
     \end{tabular}};
  \draw[rulegrey, line width=0.4pt] (14.30,8.86) -- (17.25,8.86);
  \node[quiet, anchor=north west, align=left, text width=2.90cm] at (14.30,8.74)
    {Only the wedge areas are drawn to scale. $V_{\mathrm{pure}}$ is
     orthogonal to the context subspace, so the two wedges share no direction.
     Cell line, the positive control, lies inside it, at $15$--$22\times$ the
     slab's share.};

% ###########################################################################
% # the divider
% ###########################################################################
  \draw[rulegrey, line width=0.4pt] (0,5.05) -- (17.25,5.05);
  \node[anchor=north west, font=\scriptsize\itshape, text=quietink]
    at (0,4.99) {three intervention designs failed on this geometry, and each
                 failure is a constraint on the fourth};

% ###########################################################################
% # INSET 1 -- mismatched geometry
% ###########################################################################
  \node[ttl, anchor=north west] at (0,4.60) {ONE\quad MISMATCHED GEOMETRY};

  \begin{scope}[shift={(1.25,2.92)}]
    \slabdisc{1.05}
    % 86 deg to the slab axis: cos^2(86) = 0.0049, i.e. 0.5% of the energy in
    % the slab and 99.5% outside it. The angle IS the number.
    \draw[ctrlarm] (0,0) -- (86:1.32);
    \draw[rulegrey, line width=0.3pt, densely dotted] (86:1.32) -- (0.092,0);
    \draw[accent, line width=1.8pt] (0,0) -- (0.092,0);
    \draw[lead] (0.12,0.02) -- (1.15,0.38);
    \node[quiet, anchor=west] at (0.16,1.15) {control};
  \end{scope}

  \node[quiet, anchor=west, align=left, text width=2.45cm] at (2.45,3.30)
    {its part inside the slab: $0.5\%$ of its energy};

  \node[lab, anchor=north west, align=left, text width=5.20cm] at (0,1.92)
    {The control was drawn in the full $2048$-dimensional stream while the
     intervention acts in at most ten directions, so $\approx 99.5\%$ of its
     energy lay where the hook never reached. ``The swap moves the prediction
     less than noise'' is what that geometry predicts, whatever the readout
     does.};

% ###########################################################################
% # INSET 2 -- ablate-then-add
% ###########################################################################
  \node[ttl, anchor=north west] at (6.00,4.60) {TWO\quad ABLATE-THEN-ADD};

  \begin{scope}[shift={(7.25,2.92)}]
    \slabdisc{1.05}
    \draw[lead] (\slabHalf:1.05)  -- (1.55,1.02);
    \draw[lead] (-\slabHalf:1.05) -- (1.55,-0.68);
  \end{scope}

  \begin{scope}[shift={(8.85,2.85)}]
    % added: an uncentred class mean = the shared global mean + its own part
    \node[quiet, anchor=south west] at (0,1.02) {added: $\mu_B$, uncentred};
    \fill[rulegrey!45] (0,0.84)    rectangle (1.25,1.00);
    \fill[accent]      (1.25,0.84) rectangle (2.00,1.00);
    % removed: the cell's own slab component, carrying the same global mean
    \node[quiet, anchor=south west] at (0,0.48) {removed: $\Pi h$};
    \fill[rulegrey!45] (0,0.30)    rectangle (1.25,0.46);
    \fill[accent]      (1.25,0.30) rectangle (1.70,0.46);
    \draw[rulegrey, line width=0.4pt, decorate,
          decoration={brace, amplitude=3pt, mirror}] (0,0.24) -- (1.25,0.24);
    \node[quiet, anchor=north] at (0.62,0.08) {$\bar m$, shared};
    % what the hook actually delivers
    \node[quiet, anchor=south west] at (0,-0.62) {what the hook delivers};
    \fill[accent] (0,-0.80) rectangle (0.30,-0.64);
    \node[quiet, anchor=west] at (0.40,-0.72) {$\bar m$ cancels};
  \end{scope}

  \node[lab, anchor=north west, align=left, text width=5.20cm] at (6.00,1.92)
    {The hook ablates before it adds, so what is delivered is
     (added $-$ removed), and an uncentred class mean shares its global-mean
     component with the thing removed. Substitution becomes restoration, and
     the control arm --- a real displacement --- delivers $5.3\times$ the
     perturbation.};

% ###########################################################################
% # INSET 3 -- the normaliser
% ###########################################################################
  \node[ttl, anchor=north west] at (12.00,4.60) {THREE\quad THE NORMALISER};

  \begin{scope}[shift={(13.25,2.92)}]
    \slabdisc{1.05}
    \draw[ctrlarm] (0,0) -- (2.2:0.95);
    \draw[lead] (\slabHalf:1.05)  -- (1.20,0.84);
    \draw[lead] (-\slabHalf:1.05) -- (1.20,0.06);
  \end{scope}

  % The four targets are drawn identically on purpose: the normaliser does not
  % know which of them carries identity, and that is the whole failure.
  \begin{scope}[shift={(14.45,2.92)}]
    \node[quiet, anchor=south west] at (0,0.84) {clean $\log P$};
    \draw[rulegrey, line width=0.4pt] (0,0.78) -- (1.70,0.78);
    \draw[rulegrey, line width=0.4pt] (0,0.06) -- (1.70,0.06);
    \foreach \x in {0.20,0.65,1.10,1.55}{
      \fill[ink] (\x,0.78) circle (1.5pt);
      \fill[ink] (\x,0.06) circle (1.5pt);
      \draw[quietink, line width=0.5pt, -{Latex[length=2.6pt,width=2.2pt]}]
        (\x,0.72) -- (\x,0.14);}
    \node[quiet, anchor=north west] at (0,-0.06) {the control arm};
    % anchor=west under rotate=90 puts the START of the string at this point and
    % grows it UPWARD, so the label spans the two rules it measures between
    % instead of hanging below the lower one. Centred at 0.42 it reached y=-0.23
    % and collided with the descender line of "the control arm".
    \node[quiet, rotate=90, anchor=west] at (2.04,0.10) {$+2.75$ nats};
  \end{scope}

  \node[lab, anchor=north west, align=left, text width=5.20cm] at (12.00,1.92)
    {A random direction \emph{inside} the slab inflates the logsumexp
     normaliser by a measured $+2.75$ nats, where the observed class-mean
     displacement does not. Every target is depressed by it, carrying identity
     or not: norm-matched is not disruption-matched.};

  % bounding box, so the picture is exactly the fullwidth measure
  \path (0,0.05) rectangle (17.28,11.75);
\end{tikzpicture}

\caption[The drug slab and the context subspace, by variance
share]{\textbf{The purified drug slab is a thin wedge of the residual stream
beside the context subspace, and each of the three failed identity tests turns
on that geometry.} The disc is the
$2048$-dimensional residual stream at one layer. \emph{The two wedge areas are
the only scaled marks in the figure}, drawn to the measured subspace variance
shares in \texttt{workspace\_probe.json} (twelve drug classes, forty cells
each, chance $0.083$), averaged over the seven probed layers: $0.766$ for the
context subspace of cell line, plate and dose, and $0.031$ for the purified
drug slab $V_{\mathrm{pure}}$. Cell line alone, the positive control, reads
$0.610$ and lies inside the context subspace, so it is quoted as a number
rather than drawn as a third wedge. $V_{\mathrm{pure}}$ is what remains after
the drug slab --- spanned by at most eleven directions, since it is built from
twelve centred class means, of which the top ten are kept --- is orthogonalised
against that context subspace, which is why the two wedges share no direction.
The unshaded remainder closes the disc and carries no number. Everything else
is schematic and has no scale: the slab is $11.3^{\circ}$ wide and cannot hold
three class means at true size, so its interior is magnified, where the
\textsc{swap} and \textsc{permuted} displacements are drawn at equal length
because equal length is what makes \textsc{permuted} a matched control. The one
exception among the insets is inset one's angle, which is drawn true: a
direction drawn isotropically in $2048$ dimensions puts $10/2048$ of its energy
in a ten-dimensional slab, the appendix's $\approx 99.5\%$ lying outside it. No
decode accuracy, layer axis or KL divergence appears here; those are
\cref{fig:mechanism}. The three insets are the failed designs of
\cref{sec:app-identity}.}
\label{fig:slab}
\end{fullwidth}
\end{figure}



\paragraph{Failed design one, mismatched geometry} The first test overwrote drug
A's code with drug B's inside a $\leq 10$-dimensional purified subspace, and
compared the resulting KL against a random vector of identical norm drawn in the
\emph{full} $2048$-dimensional stream. Roughly $99.5\%$ of the control's energy
then lay where the intervention never reached, including the cell-line
directions at $\approx 0.6$ variance share. ``The swap moves the prediction less
than noise'' is what that geometry predicts, whatever the readout does.

\paragraph{Failed design two, ablate-then-add} The second drew the control
inside the subspace, which looked like a repair. The hook ablates
before adding, so what is delivered is (added $-$ removed), and an uncentred
class mean shares its global-mean component with the thing removed. The
substitution arm was then close to a restoration, the control arm a real
displacement. Simulated with no model present, the control delivers $5.3\times$
the perturbation in $100\%$ of draws, predicted gap
$2\langle \mu_B, u\rangle = 9.179$ against an observed $9.162$, with $\mu_B$ the
injected class mean and $u$ the cell's component in the slab. The retracted
result was geometry, not biology.\corrmark{The ablate-then-add identity result,
retracted as an artefact of the hook's own algebra.}

\paragraph{Failed design three, the normaliser} The third injected the
displacement and scored $\Delta \log P(g_B)$ against a norm-matched random
direction, and was confounded the opposite way. A random direction inside the
slab inflates the logsumexp normaliser (measured $+2.75$ nats) where the
observed class-mean displacement does not, so every target is depressed under
the control arm whether or not it carries identity. With the scored target
stripped of all identity information, $36\%$ of the effect still stood.
Norm-matched is not disruption-matched.

That argument also retires the simplest \emph{injection} test, ``how far did the
prediction move''. Distance moved is what norm-matching equalises, so the test
has no power in precisely the world where the readout does read the drug. It
does not reach the behavioural output-invariance instrument of
\cref{sec:methods-controls}, which perturbs the prompt and matches no norms;
that instrument's own limitation, registering whether the generated text moves
and not whether it moves correctly, is stated where it is used.

\paragraph{Three further guards are engineering, not inference} A
\emph{drug-proxy screen} refuses any context label that is a near-bijection with
drug; it caught \texttt{dose\_float} holding a sample identifier in the residual
builders, a label that would have deleted the drug axis had the builders
orthogonalised against it. An \emph{injected-norm diagnostic} reports the
displacement against the cell's centred slab component ($0.69$--$1.40$), so the
injection is the size of the thing it displaces and a null cannot be blamed on a
whisper. A \emph{selftest with teeth} runs three planted worlds
(drug inert, drug read, and drug inert behind a readout sensitive to magnitude
and not direction) through the headline code path, separates them ($-0.042$,
$+1.520$, $-0.010$), and failed four times in development, each failure a real
defect the guard list now covers.

\begin{scopelimit}[carried into \cref{sec:res-mechanism}]
The third planted world is not a normaliser control. Each planted response is a
single token, so nothing diverges, and the state-dependent part of the
normaliser is precisely what the selftest cannot exercise. What separating the
three worlds licenses is the code path the selftest runs them through.
\end{scopelimit}

% ===========================================================================
\framedecl{ER}
\section{The metric exploit is a property of the metric, not of one baseline}
\label{sec:app-selftest}

\begin{question}
A metric audit whose own instrument is untested proves nothing, and the exploit
it turns on was demonstrated with a single construction. The instrument here is
a second, independent zero-information baseline put through the same framework.
It saturates the standard metric too.
\end{question}

\noindent
The self-test does the deterministic things first (a perfect prediction reaches
the metric's maximum, a zero shift and a flat truth return no score), then the
thing that matters. It asserts the exploit under a \emph{second, independent}
zero-information baseline, where a leave-one-out mean scores $\DEdr = 1.000$
while $\NIR$ refuses to place it above chance. The coded assertion is
$\NIR \leq 0.6$; the recorded value is $0.000$.

\paragraph{What that $0.000$ is not} It is not the metric's chance value,
$0.500$ (\cref{sec:methods-metrics}), which is what the generic arm scores in
\cref{tab:lookup} and \cref{tab:allarms}. Nor is it the untied degeneracy that
section describes, since the self-test uses the same half-tie convention as
everything else. It is a property of this construction. The mean of the five
other drugs resembles those five more closely than the held-out one, so it loses
every comparison.

Two unrelated constructions saturate the standard metric, and the discrimination
metric places neither above chance. One would leave the standard metric merely
unlucky; two make the exploit a property of the metric.
